\documentclass{article}
\usepackage[utf8]{inputenc}
\usepackage[margin=1in]{geometry}
\usepackage{graphicx}
\usepackage{amsmath}
\usepackage{amsfonts}
\usepackage{amssymb}
\usepackage{algorithm}
\usepackage{algorithmic}
\usepackage{booktabs}
\usepackage{url}
\usepackage{fancyhdr}
\usepackage{hyperref}
\usepackage{subcaption}
\usepackage{xcolor}
\usepackage{enumitem}

\pagestyle{fancy}
\fancyhf{}
\rhead{Social Deduction Game Manual}
\lhead{\leftmark}
\cfoot{\thepage}

\begin{filecontents}{references.bib}
@inproceedings{werewolf2009,
  title={One Night Ultimate Werewolf: A fast social deduction game},
  author={Alspach, Ted},
  booktitle={BoardGameGeek},
  year={2009}
}

@article{mafia1986,
  title={Mafia: A party game},
  author={Plotkin, Dimitry Davidoff},
  journal={Psychology Department, Moscow State University},
  year={1986}
}

@inproceedings{resistance2009,
  title={The Resistance: A social deduction board game},
  author={Eskridge, Donald X},
  booktitle={Indie Boards and Cards},
  year={2009}
}
\end{filecontents}

\title{[TITLE TO BE FILLED]}
\author{AI Game Designer}
\date{\today}

\begin{document}

\maketitle

\begin{abstract}
[ABSTRACT TO BE FILLED]
This manual presents a comprehensive guide to playing [GAME NAME], a novel social deduction game that combines traditional deductive reasoning with innovative mechanics. The game supports [X] players and features unique roles, strategic depth, and engaging social interaction.
\end{abstract}

\section{Introduction}
[INTRODUCTION TO BE FILLED]

Social deduction games have captivated players for decades, from the classic Mafia \cite{mafia1986} to modern variants like One Night Ultimate Werewolf \cite{werewolf2009} and The Resistance \cite{resistance2009}. This manual introduces [GAME NAME], a new social deduction game that brings fresh mechanics and engaging gameplay to the genre.

\section{Game Overview}

\subsection{Players and Setup}
[SETUP DETAILS TO BE FILLED]

\subsection{Objective}
[OBJECTIVE TO BE FILLED]

\section{Components}
[COMPONENTS TO BE FILLED]

The game includes:
\begin{itemize}
\item Role cards for each player type
\item Reference cards explaining abilities
\item Game master instructions
\item Phase tracking tokens (if applicable)
\end{itemize}

\section{Roles and Abilities}
[ROLES TO BE FILLED]

This section details each role in the game, including their unique abilities, objectives, and strategic considerations.

\subsection{[ROLE NAME]}
[ROLE DESCRIPTION TO BE FILLED]

\section{Game Flow}
[GAME FLOW TO BE FILLED]

\subsection{Setup Phase}
[SETUP INSTRUCTIONS TO BE FILLED]

\subsection{Game Phases}
[PHASE DESCRIPTIONS TO BE FILLED]

\section{Rules and Mechanics}
[RULES TO BE FILLED]

\subsection{Discussion Phase}
[DISCUSSION RULES TO BE FILLED]

\subsection{Voting}
[VOTING RULES TO BE FILLED]

\subsection{Special Abilities}
[ABILITY RULES TO BE FILLED]

\section{Victory Conditions}
[VICTORY CONDITIONS TO BE FILLED]

\section{Strategy and Tips}
[STRATEGY SECTION TO BE FILLED]

\subsection{General Strategy}
[GENERAL TIPS TO BE FILLED]

\subsection{Role-Specific Tips}
[ROLE-SPECIFIC STRATEGY TO BE FILLED]

\section{Variants and Optional Rules}
[VARIANTS TO BE FILLED]

\section{Designer Notes}
[DESIGNER NOTES TO BE FILLED]

This game was designed using AI-assisted game design principles, iteratively tested and refined to create an engaging social deduction experience.

\section{Conclusion}
[CONCLUSION TO BE FILLED]

\bibliographystyle{plain}
\bibliography{references}

\end{document} 